\documentclass[10pt,a4paper]{amsart}
\usepackage[margin=19mm]{geometry}
\usepackage{amsmath,amssymb,mathtools}
\usepackage[scr=rsfs,cal=euler]{mathalfa}
\usepackage{xfrac}
\usepackage[unicode,hidelinks,bookmarks=false]{hyperref}
\newcommand{\ie}{i.e.\ }
\newcommand{\cf}{cf.\ }
\newcommand{\resp}{respectively\ }
\newcommand{\Subsection}{\subsection}
\providecommand{\nobreakdash}{\nobreak\hbox{-}}
\setlength{\emergencystretch}{2em}
\multlinegap0pt
\usepackage{bm}
\usepackage{bbm}
\usepackage{dsfont}
\usepackage{xspace}
\usepackage{cancel}
\usepackage{mathtools}
\usepackage{amsthm}
\makeatletter
\@ifundefined{lemma}{\newtheorem{lemma}{Lemma}}{}
\@ifundefined{theorem}{\newtheorem{theorem}[lemma]{Theorem}}{}
\makeatother
\newcommand{\dd}{\mathrm{d}}
\title[Hybrid Risk Processes: A Versatile Framework for Modern Ruin]{Hybrid Risk Processes: A Versatile Framework for Modern Ruin Problems}
\author{Oscar Peralta, Habacuq Vallejo}
\date{Source: arXiv:2506.23685v1 (CC BY 4.0)}
\begin{document}
\maketitle

\noindent\emph{Source and licence.} Hybrid Risk Processes: A Versatile Framework for Modern Ruin Problems. Version arXiv:2506.23685v1 (CC BY 4.0), distributed under the Creative Commons Attribution 4.0 International licence, as recorded in the arXiv licence metadata for this exact version and shown on the abstract page https://arxiv.org/abs/2506.23685v1. Full rights notice: licenses/arxiv-2506.23685-CC-BY-4.0.txt.\par\smallskip

\noindent\textit{Editorial scope.} Counted unit: the Theorem whose source label is \texttt{thm:jump\_density\_restructured} in the article (source0.tex lines 247--255, source label \texttt{thm:jump\_density\_restructured}), delivered together with the article's own complete original proof of it (source0.tex lines 256--272, one \texttt{proof} environment of 17 lines). No mathematics is written, repaired or re-derived: statement and proof are the article's own text line for line. Required context retained uncounted: none. Every definition, display and label of those lines is retained unchanged. Omitted from this package and not redistributed: 1--246, 273--704. Nothing inside a retained excerpt is omitted or altered: every retained body is delivered exactly as the article's lines are written, so no omission receipt of any kind accompanies this package. Citations: the retained excerpts carry no citation command at all, so no bibliography entry of the article is transcribed and no reference is redistributed. Reference adaptation: none was needed and none was made; every reference inside the delivered unit resolves inside it, so no unresolved reference or citation is produced. Printed slips kept literal: no printed word, symbol, hypothesis or number of the counted statement or of its proof was altered. Layout and class: the article's own class is \texttt{article} with packages including amsmath, epsfig, graphicx, color, float, amssymb, amsthm, subcaption, relsize, comment, hyperref, pdfpages, mwe, algorithm; the wrapper is the lane's pinned \texttt{amsart} editorial wrapper at 10pt on A4, which restates the theorem-like environments the retained text needs and adds the article's own macro definitions that the retained text needs (\texttt{\textbackslash dd}), with extra wrapper packages (bm, bbm, dsfont, xspace, cancel, mathtools). The delivered numbering is the unit's own. Assets: no retained span contains a figure environment, a graphics-inclusion command, a PDF-page inclusion, a table or a TikZ picture; no image file is copied into this package, the package declares no asset dependency and no image file is redistributed with it. Licence: version arXiv:2506.23685v1 of this article is distributed under the Creative Commons Attribution 4.0 International licence, as recorded in the arXiv licence metadata for this exact version and shown on the abstract page https://arxiv.org/abs/2506.23685v1. A full rights notice is delivered at licenses/arxiv-2506.23685-CC-BY-4.0.txt. Characters: every retained excerpt is pure ASCII, so the delivered engine, XeLaTeX with its default Latin Modern text font, reports no missing character.
\begin{theorem}
\label{thm:jump_density_restructured}
For \(y > 0\) and \(i,j \in \mathcal{S}^p\), the intensity of a downward jump of size approximately \(y\) from level \(x\) is given by
\begin{align*}
&\mathbb{P}(\text{downward jump in }(t, t+\dd t)\text{ of size in }(y, y+\dd y), L_{t+\dd t}=j \mid R_t=x, L_t=i) = \\
&\quad \bm{e}_i^\intercal \bm{\Lambda}_{\mathcal{S}^p\mathcal{S}^-}(x) \left( \prod_{u=0}^y (\bm{I} + \bm{\Delta}_\mu(x-u)\bm{\Lambda}_{\mathcal{S}^-\mathcal{S}^-}(x-u)\dd u) \right) \bm{\Delta}_\mu(x-y)\bm{\Lambda}_{\mathcal{S}^-\mathcal{S}^p}(x-y) \bm{e}_j \dd y \dd t,
\end{align*}
where \(\bm{\Delta}_\mu(z) = \text{diag}(|1/\mu(k,z)|\,:\,k\in \mathcal{S}^-)\) is the diagonal matrix of inverse drift rates for each state \(k \in \mathcal{S}^-\) at level \(z\). An analogous result holds for upward jumps.
\end{theorem}

\begin{proof}
We provide the proof for the downward jump case; the argument for an upward jump is analogous. A downward jump in \((L,R)\) corresponds to an excursion of \((J,X)\) through the state space \(\mathcal{S}^-\). The process unfolds in three steps:
\begin{enumerate}
    \item The process \(J\) transitions from a state \(i \in \mathcal{S}^p\) to an initial state within \(\mathcal{S}^-\). The matrix of intensities for this entry is \(\bm{\Lambda}_{\mathcal{S}^p\mathcal{S}^-}(x)\).
    \item The process \(J\) sojourns in \(\mathcal{S}^-\). During this sojourn, while the environment is in state \(k \in \mathcal{S}^-\), the surplus \(X\) drifts downward at a state- and level-dependent speed given by \(\mu(k,X_s)\). The total accumulated drift over the sojourn duration determines the jump size \(y\).
    \item The process \(J\) exits \(\mathcal{S}^-\) by jumping to a state \(j \in \mathcal{S}^p\).
\end{enumerate}
To find the density of the jump size \(y\), we consider the evolution of the environmental process \(J\) with respect to the accumulated jump size itself, rather than sojourn time. When the accumulated jump size is \(y\) and the environment is in state \(k\), the surplus level is \(x-y\). At this point, the speed at which the jump size increases is given by \(|\mu(k,x-y)|\). To accumulate a further infinitesimal jump of size \(\dd y\), it takes an infinitesimal amount of sojourn time \(\dd s = \dd y / |\mu(k,x-y)|\).

The intensity matrix for transitions within \(\mathcal{S}^-\) with respect to sojourn time \(s\) is \(\bm{\Lambda}_{\mathcal{S}^-\mathcal{S}^-}(x-y)\). To find the intensity matrix with respect to the jump size \(y\), we scale the time-based intensities by the amount of time required to accumulate a unit of jump size. This yields an intensity matrix with respect to \(y\), given by the product
\[
\bm{T}^-(y;x) = \bm{\Delta}_\mu(x-y) \bm{\Lambda}_{\mathcal{S}^-\mathcal{S}^-}(x-y).
\]
The matrix of transition probabilities for the environment over a jump of size \(y\) is then given by the product integral \(\prod_{u=0}^y (\bm{I} + \bm{T}^-(u;x)\dd u)\).

Similarly, the matrix of exit intensities from \(\mathcal{S}^-\) to \(\mathcal{S}^p\) per unit of jump size is \(\bm{\Delta}_\mu(x-y)\bm{\Lambda}_{\mathcal{S}^-\mathcal{S}^p}(x-y)\). Combining the intensity of entry, the propagation during the sojourn (the product integral), and the intensity of exit gives the overall intensity stated in the theorem.
\end{proof}

\end{document}
